\section{Related Work}
\label{sec:related}

\paragraph{\it Underapproximate Reasoning} A number of recently
proposed logical frameworks provide program reasoning principles based
on under-, rather than over-, approximate abstractions.  Incorrectness
Logic~\cite{OHP19} and related variants~\cite{LR+22,RBD+20,ZDS23}
present compositional proof rules that enable verification of
reachability (aka incorrectness) properties of first-order imperative
programs; their primary motivation is to enable formal reasoning of
(unsound) program analysis tools such as Infer~\cite{DFLO19}.  Ideas
related to these notions have also been incorporated into type
systems~\cite{RW24,CoverageType,JY15,LPR26} for functional languages,
but none of these efforts have considered their application to the
synthesis of effectful test generators of the kind proposed in this
paper.  For example~\cite{CoverageType} uses underapproximation to
typecheck that a PBT generator is \emph{complete} with respect to the
functional inputs it generates to the SUT, while ~\cite{LZ+25}
considers an enumerative repair strategy to patch incomplete
generators so that they become complete.  In contrast, in this paper,
we show how interpreting type specifications from the lens of
underapproximation can guide a synthesis procedure to construct
bespoke effectful test generators tailored to the property of
interest, informed by the behavior of the operators provided by the
SUT.  As a mechanism to characterize underapproximate behavior, types
identify \emph{sufficient} effectful constraints that any test input
sequence produced by the generator \emph{must} satisfy based on these
specifications.

\paragraph{{\it Specializing PBT Generators}}
There has been significant prior work to make PBT implementations more
effective, by providing tools that enable reasoning and customization
of effect traces produced by a test generator.  For example,
~\citet{CP+09} describes a customizable user-level scheduler that uses
Quviq QuickCheck's support for state machine models to implement
deterministic replayable test inputs to an Erlang SUT.
Concuerror~\cite{CGS13} systematically explores process interleavings
in a concurrent Erlang program via program instrumentation and
preemption bounding.  Quickstrom~\cite{Quickstrom} uses specifications
of possible actions expressed in its Specstrom specification language,
a dialect of LTL, to dynamically choose the next input that the
generator should provide for interactive applications.  In contrast,
\name's type-guided synthesis strategy yields
generators influenced by both the global property of interest and the
local \tyName{} specifications of the operations used by the SUT.
Importantly, the synthesis procedure is guided by the history and
future SREs of provided specifications to determine a semantically
meaningful ordering of effectful actions.

%% TLA+~\cite{Lam02} is a specification language based on LTL for
%% modeling finite-state distributed systems; the correctness of these
%% specifications are verified using the TLC explicit-state model
%% checker. TLA+ and its associated tooling has had notable real-world
%% impact~\cite{NR+15}.  While \name{}'s use of LTL$_f$ specifications in
%% \PAT{}s is a point of commonality with TLA+, the integration of these
%% specifications within a refinement type system, their role in driving
%% a component-based synthesis procedure, and the top-down (TLA) vs.
%% bottom-up (\name{}) exploration mechanism, differentiates \name{}'s
%% motivation and design from TLA+ and TLC in obvious ways.


\paragraph{Type and Effect Systems}
Type and effect systems that target \emph{temporal} properties on the
sequences of effects that a program may \emph{produce} is a
well-studied subject. For example, \citet{SS+04} presents a type and
effect system for reasoning about the shape of histories (i.e., finite
traces) of events embedded in a program.  \citet{KT+14} present a type
and effect system that additionally supports verification properties
of infinite traces, specified as B\"{u}chi automata.  More recently,
\citet{Temporal-Verification} have considered how to support richer
control flow structures, e.g., delimited continuations, in such an
effect system.  Closest to our work are \emph{Hoare Automata Types}
(\textsc{HAT}s)~\cite{ZYDJ24}, which integrate symbolic finite
automata into a refinement type system. \textsc{HAT}s enable reasoning
about stateful sequential programs structured as a functional core
interacting with opaque effectful libraries.  \tyName{} revise
\textsc{HAT}s by viewing them as underapproximate specifications -
every trace accepted by the future regex in a \tyName{} is
justified by a trace accepted by its history.


%% \paragraph{{\it Verification}} Formally proving the correctness of
%% distributed protocols and models has long been a topic of significant
%% interest~\cite{KA+07,HH+15,SWT18}.  These approaches provide strong
%% correctness guarantees at the cost of significant investment on the
%% part of the proof engineer, who is responsible for, e.g., defining
%% suitable inductive invariants for the verification
%% task~\cite{PM+16,YR+21,MG+19}.  In contrast, our focus in this work is
%% to improve the effectiveness of falsification techniques--- validating
%% the presence of bugs in a distributed protocol design, rather than
%% their absence. In this sense, we are more closely related to recently
%% proposed approaches for formally reasoning about
%% incorrectness~\cite{OHP19, LR+22, RBD+20, casl}.  While
%% \name{} cannot verify the correctness of a model, the burden we impose
%% on test engineers, i.e., providing handler specifications as PAT{}s,
%% as well as a global safety/liveness property in LTL$_f$, is
%% significantly less than what is required to verify full functional
%% correctness of these designs.

%% \paragraph{{\it Testing}} Outside of the aforementioned P
%% language~\cite{DGJ+13,ModP}, several other efforts have considered how
%% to improve the capabilities of testing frameworks for distributed
%% systems.  Jepsen~\cite{Jepsen} is a randomized testing system that
%% seeks to reveal bugs when an application is deployed on a
%% weakly-consistent storage system; \citet{OM+18} defines a randomized
%% testing procedure for message-passing distributed systems with
%% guaranteed lower bounds on the probability of finding a depth-$d$ bug,
%% where $d$ is the minimum length of the sequence of events sufficent to
%% witness the error.  Morpheus~\cite{YY20} uses partial order sampling
%% and conflict analysis to control scheduling
%% decisions. MonkeyDB~\cite{RK+19} uses a demonic scheduling mechanism
%% to expose safety violations in SQL applications that interact with a
%% weakly-isolated storage backend.  Clotho~\cite{RK+19} combines static
%% analysis with a bounded model-checker to generate tests that expose
%% serializability violations in weakly-consistent database systems.
%% While these efforts are all agnostic to the property under test,
%% \name{}'s property-guided synthesis procedure derives a controller
%% specialized to the target property and handler specifications that
%% capture temporal dependencies between actors. In this sense, our
%% approach can be seen as a form of property-based testing
%% (PBT)~\cite{CH00, HC+24} applied to open reactive systems.  Broadly
%% related to our approach is Mocket~\cite{WD+23}, a PBT-style testing
%% framework that uses the state space graph extracted from
%% model-checking TLA+ specifications~\cite{Lam02} to force executions to
%% follow specific paths in the graph.  Unlike \name{}, Mocket requires
%% manual instrumentation of implementations to align actions defined in
%% the specification with the corresponding code in the implementation,
%% and relies on the TLC model-checker to produce the state space graph.
%% In constrast, \name{} uses a compositional refinement type system to
%% drive synthesis, and requires no instrumentation or \emph{a priori}
%% enumeration of the state space to synthesize its controllers.


\section{Conclusions}\label{sec:conc}

This paper proposes trace-guided synthesis of effectful PBT
generators.  Our key innovations are (i) the adaption of symbolic
regular expressions, expressed as types, that define traces of
effectful actions that occur before and after the execution of an
expression, as underapproximate specifications, and (ii) the synthesis
of bespoke test generators derived from these specifications.
Experimental results on a wide range of sequential, concurrent, and
distributed applications, including real-world system models used in
production, show that \name\ is significantly more effective in
identifying test input sequences that falsify a property than the
existing state-of-the-art.
